\section{The source-defined arithmetic candidate and its inputs}\label{sec:induction}
\subsection{Do not identify two candidates by notation}
The direct matrix recursion and the classical arithmetic-induction route start with different data. The first produces an actual type satisfying a restricted construction contract. The second starts from a complete source character and uses classical induction and descent to produce a candidate. We denote the latter by $\pi_{\mathrm{BH}}(\rho)$ in this section. This notation does not identify it with any independently constructed compact type.

The audited source-defined branch is over $\F_3((t))$ and uses the original rank-nine source. It is distinct from the characteristic-zero continuation113 candidate discussed in Part II. Its existence through classical inputs does not prove the corrected-character occurrence question for that other candidate, nor does it place a prescribed native packet in the current one.

\subsection{The degree-nine field does not initially have a cubic tower}
Write $A=t^{-1}$, and retain the projective coordinates
\[
 x^9+x=A,\qquad z^9+z=A^2.
\]
The source record takes $y=z-x-x^3$. In characteristic three,
\[
 y^9+y=B,\qquad B=A^2-A-A^3.
 \tag{25.1}\label{eq:B}
\]
The inducing field is $M=F(y)$, of degree nine. Its retained Galois-closure analysis has translation subgroup $(C_3)^2$ with a $C_4$ action irreducible over $\F_3$. There is no cubic intermediate field over $F$. Therefore one cannot simply write two cubic automorphic-induction arrows over $F$ without changing the base and proving the resulting descent.

Pass to the unramified extension $K=\F_{81}((t))$. The constants now split the translation action. Choose $a\ne0$ with $a^9=-a$ and $i$ with $i^2=-1$, and define
\[
 r=\frac{y^3-a^2y}{ia^3},\qquad s=\frac ya.
\]
Then
\[
 r^3-r=\frac{B}{ia},\qquad s^3-s=ir.
 \tag{25.2}\label{eq:tower}
\]
To check the first equality, use $a^8=-1$ to cancel the $y^3$ terms after cubing. The second follows immediately from the definition of $r$. These identities exhibit the cyclic cubic tower over $K$ used by the classical route. The degree and regularity assertions remain part of the complete source fixture.

The inducing character is also complete data. The inherited record includes its ramified Artin--Schreier class, its unit values, and an actual unramified factor. Replacing it by the displayed pole order alone loses information needed for the resulting central character and noncompact extension. The noted different exponent $24$ and character conductor $9$ give the expected induced ordinary conductor $33$, but do not determine that full character.

\subsection{Induction, invariance, and descent are separate gates}
The construction first applies cyclic automorphic induction at the two cubic stages over $K$. Cuspidality requires the recorded regularity at each relevant step. It then invokes the classical independence of the resulting monomial construction from an admissible cyclic tower. This independence is a substantial theorem. It is what allows a Galois conjugate of the chosen tower to be compared with the original construction. The fact that the source descends to $F$ does not, by itself, make each arbitrarily chosen tower Frobenius invariant.

Only after this invariance has been established does the tame descent theorem apply. Its fibre is described by the unramified characters associated to $K/F$, here of order dividing four. The specified determinant of $\rho$ selects the desired central character.

\begin{lemma}[Unique fixed-center member in a coprime descent fibre]\label{lem:descent}
Let a descent fibre be a torsor for a cyclic character group $X$ of order $m$. Suppose twisting a rank-$n$ member by $\alpha\in X$ changes its central character by $\alpha^n$. If $\gcd(n,m)=1$, then any attainable target central character selects exactly one member of the fibre.
\end{lemma}
\begin{proof}
The map $\alpha\mapsto\alpha^n$ is a bijection of $X$ because multiplication by $n$ is invertible modulo $m$. Given one member and an attainable ratio of central characters, there is a unique twist realizing that ratio.
\end{proof}

For $n=9$ and $m=4$, the required power map is invertible. This explains the central normalization rather than silently selecting a preferred descent. It differs from the fixed-center cubic ambiguity in continuation113, where the relevant power map is not coprime. The new finite tests include both the unique and the nonunique cases.

\subsection{What the primary-source audit does and does not establish}
The selected audit uses three layers of classical material.

\textbf{The 1999 monomial and tame-lifting statements.} In [BH99], the definitions and properties of tame lifting in Section 1.3 identify the descent fibre, twisting and central-character behavior. Section 2.3 gives the characteristic-zero monomial construction, including standard epsilon factors and pair $L$-functions and conductors. Its proof uses independence of the cyclic construction. Section 1.7 supplies the tame transfer rule for the relevant pair-conductor exponent. These statements were read at their exact pages.

\textbf{The positive-characteristic extension.} The original Section 2.3 is explicitly characteristic zero. The characteristic note on page eight of [BH13] states that the relevant earlier results extend to positive characteristic through the base-change and automorphic-induction work of [HL11]. That note was checked directly. The present audit therefore retains the positive-characteristic induction theorem as a named external input, rather than erasing the original hypothesis.

\textbf{The limit of fresh source access.} The complete [HL11] PDF could not be retrieved in this block. Its bibliographic record was accessible, but repeated full-text requests failed. The detailed statements III.2 and III.3, including the restricted pair statement used in [Q15], retain their inherited source-audit status. They are not labeled freshly reread here. This is a narrower limitation than saying the entire arithmetic route was unaudited.

There are two particularly important exclusions. First, the effective construction in [BH13] treats the totally wildly ramified correspondence as given at its starting point. That paper's effective theorem is not used here as an independent solution of the same totally wild input. Its characteristic note is the relevant source interface. Second, the equality of pair conductors in [BH99, 2.3(vii)] is not equality of full pair epsilon factors. The introduction expressly distinguishes these strengths.

\subsection{The resulting candidate still has a native attachment obligation}
Under the recorded source and classical hypotheses, [Q15] obtains a cuspidal candidate with prescribed central character, torsion number one, standard Swan conductor $24$, ordinary conductor $33$, and ordinary self-pair conductor $232$. The characterizing information required by the paper's H5 contract is not replaced by that list.

In particular, one cannot compare an essentially discrete source with a reducible principal series outside the scope of a pair theorem and then divide a pair conductor by nine to manufacture the standard conductor. Nor can a first-order calculation of a parameter on deep field units establish occurrence of a full simple character. The original unit domain and noncompact extension remain part of that question.

The next useful arithmetic step is therefore not another calculation of the same three integers. It is to construct or identify the actual native type and its complete character in this fixed candidate, and to compare its marked data with the supplied source. That is the boundary the finite results in the following sections are designed to expose.

\source{[Q15] is the inherited source-defined constructor audit. [BH99, printed pp. 15, 18, 23--25, 53--54] and [BH13, p. 8] were freshly checked. [HL11] remains a named classical input with an explicit fresh-access limitation. No source-defined candidate is substituted for a separately fixed type candidate without a proof of their identification.}
